\section{Pruebas y CI/CD}
\label{sec:pruebas-cicd}

\subsection{Pirámide de pruebas}

La estrategia de pruebas cubre las cuatro capas clásicas de la pirámide~\cite{ammann2016testing},
más una capa de carga, cada una respondiendo una pregunta que las demás no pueden:

\begin{table}[H]
\centering
\small
\caption{Pirámide de pruebas del sistema, con resultados reales.}
\label{tab:piramide}
\begin{tabular}{@{}p{2.3cm}p{3.4cm}p{4.2cm}p{3.2cm}@{}}
\toprule
\textbf{Capa} & \textbf{Herramienta} & \textbf{Pregunta que responde} & \textbf{Resultado} \\
\midrule
Unitaria & JUnit 5 / Vitest / JUnit Android & ¿Esta clase, aislada, hace lo que dice? & 90/90 backend (5 servicios Java), 78/78 web, 5/5 móvil \\
Integración & Testcontainers & ¿Funciona contra un CockroachDB real, no un \emph{mock}? & \texttt{TicketRepositoryIntegrationTest} \\
Contrato & Pact (consumidor + proveedor) & ¿El backend sigue la forma exacta que el cliente espera? & 2/2 consumidor, 2/2 proveedor \\
E2E & Playwright / Compose Testing & ¿Todo el camino, navegador real incluido, funciona? & 9/9 en Chromium+Firefox+WebKit \\
Carga & Locust & ¿Aguanta concurrencia real? & Sección~\ref{sec:iso25010} \\
\bottomrule
\end{tabular}
\end{table}

\subsubsection{Pruebas de contrato}

Se adoptó \emph{consumer-driven contract testing} con Pact entre \texttt{apps/web} (consumidor) y
\texttt{ticket-service} (proveedor). El consumidor genera un archivo de contrato
(\texttt{pacts/soporte-web-ticket-service.json}) a partir de las expectativas reales del cliente
sobre la forma de la respuesta; el proveedor lo verifica reproduciendo cada interacción contra el
servicio real. Una particularidad de esta arquitectura obligó a una decisión de diseño: como
\texttt{ticket-service} valida cada JWT llamando por HTTP a \texttt{auth-service}
(\texttt{AuthGatewayFilter}) en vez de verificar la firma localmente, la verificación del
proveedor no puede correr en un contexto de prueba aislado —necesita el \emph{stack} completo
levantado y un JWT real obtenido de un inicio de sesión real.

Un segundo ajuste, encontrado al ejecutar la verificación por primera vez, fue de modelado del
contrato: los campos \texttt{technicianId} y \texttt{resolvedAt} son legítimamente nulos o no
nulos según el estado real del ticket, y un solo ejemplo con matiz de tipo (\texttt{like(null)})
no puede expresar ambos casos a la vez sobre el mismo campo. Se adoptó el patrón oficial de Pact
para colecciones con formas mixtas (\texttt{arrayContaining} con dos variantes representativas),
en vez de forzar una condición más débil o más fuerte que la que el sistema realmente cumple.

\subsection{Pipeline de CI/CD}

El archivo \texttt{.github/workflows/ci-cd.yml} define siete \emph{jobs} en un grafo dirigido
acíclico, siguiendo la cultura DevOps de que ningún cambio llegue a un artefacto publicado sin
pasar por los mismos controles~\cite{kim2021devops,humble2010continuous}:

\begin{table}[H]
\centering
\small
\caption{Los siete jobs del pipeline.}
\label{tab:cicd}
\begin{tabular}{@{}p{0.6cm}p{2.6cm}p{9.5cm}@{}}
\toprule
\textbf{\#} & \textbf{Job} & \textbf{Qué hace} \\
\midrule
1 & \texttt{lint} & ESLint (web) + Android Lint (móvil) \\
2 & \texttt{test-backend} & Los cuatro servicios Java + Node + Python, pruebas unitarias e integración \\
3 & \texttt{test-web} & Vitest, incluido el contrato Pact del consumidor \\
4 & \texttt{test-mobile} & Unitarias JVM + instrumentadas en un emulador Android headless real \\
5 & \texttt{build-images} & Construye y publica las siete imágenes Docker en GHCR (solo en \emph{push}) \\
6 & \texttt{build-mobile-apk} & Compila el APK de depuración y lo publica como artefacto \\
7 & \texttt{integration} & Levanta el \emph{stack} completo y corre Pact proveedor + Playwright contra el backend real \\
\bottomrule
\end{tabular}
\end{table}

Los \emph{jobs} 5 y 6 declaran \texttt{needs} sobre los \emph{jobs} de prueba correspondientes:
GitHub Actions no los ejecuta hasta que sus dependencias terminen en verde. Esto no siempre fue
así de completo: en una revisión externa del pipeline se encontró que \texttt{build-images} solo
dependía de \texttt{test-backend} y \texttt{test-web}, así que publicaba imágenes en GHCR aunque
\texttt{lint} o \texttt{test-mobile} estuvieran en rojo -- no era la compuerta de calidad real que
se pretendía. Se corrigió agregando esas dos dependencias faltantes. Una segunda revisión
encontró que aún faltaba la más importante: \texttt{integration} -- el único \emph{job} que
prueba contra el \emph{stack} real y no contra \emph{mocks} -- corría en \emph{paralelo} con
\texttt{build-images} en vez de antes, así que una imagen podía publicarse aunque la integración
real fallara. Se agregó también a \texttt{needs}, sin crear una dependencia circular porque
\texttt{integration} construye sus propias imágenes en local (\texttt{docker compose up
--build}), no depende de las que publica GHCR. La declaración final es
\texttt{needs: [lint, test-backend, test-web, test-mobile, integration]}, para que ningún
artefacto se publique sin que los cinco jobs de verificación pasen primero.

\subsection{Depuración real del job \texttt{integration}}
\label{sec:cicd-depuracion}

Los primeros 6 \emph{jobs} pasaban de forma consistente desde su introducción. El séptimo,
\texttt{integration}, falló en \emph{todas} las corridas sobre GitHub Actions durante varios
días, pese a que el mismo \texttt{docker compose up -d --build} que ese \emph{job} ejecuta
funcionaba sin problema en las máquinas de desarrollo del equipo. Esa asimetría -- funciona en
local, falla siempre en el \emph{runner} -- fue la pista de que la causa era estructural, no una
falla puntual. El equipo instaló la CLI oficial de GitHub (\texttt{gh}) para acceder al log real
de cada paso (la API pública solo devuelve \enquote{\texttt{Process completed with exit code 1}},
sin el comando que realmente falló), y diagnosticó tres causas raíz independientes, cada una
confirmada con evidencia real antes de corregirla:

\begin{enumerate}[nosep]
  \item \textbf{Candado circular en el healthcheck de CockroachDB.} El healthcheck de los 3
        nodos usaba \texttt{/health?ready=1} -- que solo responde \texttt{OK} una vez que el
        nodo ya pertenece a un clúster \emph{inicializado} -- mientras que \texttt{db-init} (quien
        ejecuta \texttt{cockroach init}) esperaba a que ese mismo healthcheck pasara primero.
        Nadie quedaba listo porque nadie los inicializaba, y nadie los inicializaba porque nadie
        quedaba listo. En desarrollo esto no se manifestaba porque los volúmenes de Docker ya
        traían un clúster inicializado de sesiones anteriores; un \emph{runner} de GitHub
        Actions arranca siempre desde un volumen vacío, así que el candado se activaba sin
        excepción. Un primer intento de aumentar el presupuesto de tiempo del healthcheck
        \emph{no} resolvió el problema -- solo retrasó la falla de 70\,s a 3\,min 21\,s --,
        confirmando que la causa era estructural y no de tiempo. Se corrigió cambiando el
        healthcheck a \texttt{/health} (solo \emph{liveness}) y agregando reintentos reales
        alrededor de \texttt{cockroach init}.
  \item \textbf{Datos de demo del contrato Pact, nunca versionados.} El método \texttt{@State}
        de la verificación del proveedor asumía datos de demo ya cargados por un script que
        resultó no existir en ningún commit del repositorio -- solo había existido, sin
        documentarlo como tal, en el volumen de Docker de quien lo escribió originalmente. Se
        corrigió haciendo que el propio \texttt{@State} inserte esos datos por SQL directo, que
        es precisamente para lo que existe ese mecanismo de Pact.
  \item \textbf{Playwright solo instalaba Chromium.} \texttt{playwright.config.ts} define 3
        proyectos (Chromium, Firefox, WebKit -- requisito de compatibilidad de esta entrega),
        pero el paso de instalación del \emph{workflow} tenía el argumento
        \texttt{chromium} agregado, restringiendo la instalación a un solo navegador. Se
        corrigió quitando el argumento.
\end{enumerate}

Las tres correcciones se verificaron de punta a punta contra el \emph{stack} real -- nunca
simuladas -- antes de subirlas, y el pipeline completo quedó en verde sobre un mismo commit. El
patrón común a los tres problemas es el mismo: algo que \emph{parecía} funcionar porque la
máquina de desarrollo conservaba un estado acumulado (un volumen ya inicializado, datos cargados
a mano, un navegador ya instalado) que un entorno realmente limpio nunca reproduce por su cuenta.
El detalle completo de cada diagnóstico, incluidos los registros reales y los intentos que no
funcionaron, está en \texttt{docs/ci\_cd/informe\_ci\_cd.tex}.

\subsection{Cobertura y complejidad ciclomática}

La cobertura de \texttt{ticket-service}, medida con JaCoCo, alcanza el \textbf{81.1\,\%} de
líneas del código propio, por encima del objetivo del 70\,\% que el proyecto se propuso desde la
Entrega 3. Ese número no incluye el código generado automáticamente por \texttt{protobuf}/
\texttt{grpc-java} a partir de \texttt{telemetry.proto} (\emph{getters}, \emph{setters} y
constructores mecánicos sin lógica propia, que ninguna convención de pruebas exige cubrir a
mano) — se excluye explícitamente en el \texttt{pom.xml} (perfil \texttt{coverage}), con un
comentario que documenta por qué. Sin esa exclusión, ese código generado por sí solo diluía la
cifra reportada de 59.3\,\% (antes de que \texttt{telemetry-service}/\texttt{CORREL} existieran)
a un engañoso 42.7\,\%, contando en contra código que nadie prueba por convención en ningún
proyecto Java real.

Antes de escribir pruebas nuevas para esta entrega, el código propio (sin el generado) ya estaba
en 64.4\,\%. Se identificaron con el propio reporte de JaCoCo las clases reales con mayor número
de líneas sin cubrir — no una elección arbitraria — y se agregaron 31 pruebas unitarias nuevas
(87 en total en el módulo, todas en verde) sobre, entre otras: \texttt{CorrelationService} (el
orquestador de \texttt{CORREL}, que llegó a esta entrega prácticamente sin pruebas propias), los
tres observadores concretos del patrón Observer (\texttt{EscalationLoggingObserver},
\texttt{EscalationMetricsObserver}, \texttt{EscalationEventPublisherObserver} — cerrando un hueco
real en un patrón que el propio manuscrito documenta como implementado), \texttt{
GlobalExceptionHandler}, y el mapeo/adaptador de persistencia de \texttt{Incidencia}.

La complejidad ciclomática promedio del mismo módulo, medida con la herramienta \texttt{lizard},
es de \textbf{1.8}, muy por debajo del umbral de 10 — ninguna función individual supera siquiera
el umbral de alarma de 15, resultado consistente con el refactor a métodos pequeños de la
Sección~\ref{sec:hexagonal-e4}.
